\section{Del 3: Forskjellene mellom selskapenes organisasjonskulturer}

Organisasjonskultur handler om verdier, normer, og oppfatninger som etableres gjennom erfaringer fra interne og eksterne problem, som påvirker hvordan ansatte opptrer på arbeidsplassen. Dette gjelder hvilke holdninger man har overfor arbeidet man holder på med, hvilken holdning man har til de andre ansatte, og hvordan man samarbeider med dem. Det handler om å skape rom til å utøve et yrke på best mulig måte, dette gjennom gode relasjoner og samhandling i trygge omgivelser. Organisasjoner har ulike måter for å oppnå dette på.

I denne oppgaven skal vi sammenligne to bedrifter som har forskjellig tilnærming til hjemmekontor. DNB som nylig har strammet inn på bruk av hjemmekontor, og Vend som bruker en mer fleksibel løsning for sine ansatte. Vi tar i bruk Scheins teori om organisasjonskultur i sammenligningen. Hvilke sentrale forskjeller ser man i organisasjonskultur for bedrifter som benytter seg av hjemmekontorløsningen, sett opp mot dem som ikke bruker dette?

Schein deler organisasjonskultur inn i tre ulike nivåer; artefakter, verdier og normer, og grunnleggende antakelser. Det første nivået, artefakter, er de synlige uttrykkene som man kan observere i organisasjonskulturen. I dette tilfellet er hjemmekontorordningen til de to bedriftene en slik artefakt, der DNB foretrekker fysisk oppmøte på arbeidsplassen, mens Vend har en mer fleksibel tilpasning. Verdier og normer er felles oppfatninger og uskrevne regler til hvordan de ansattes adferd forventes til være i kulturen, som også påvirker hvordan de utfører oppgaver og tar beslutninger. Innstrammingen av hjemmekontor tyder på at DNB ønsker å styrke ansattes samhandling og relasjoner på arbeidsplassen fysisk. Grunnleggende antakelser, som er dypeste og mindre synlige nivået, er organisasjonen og de ansattes virkelighetsoppfatning, altså hva de ser på som sant og usant, i tillegg til hvordan de antar ting henger sammen. Ut ifra at Vend åpner for mer hjemmekontor, kan man anta ledelsen har tillitt til at de ansatte er i stand til utføre oppgavene sine både på og utenfor arbeidsplassen.

Arbeidsplassen blir mer ettertraktet jo flere fordeler de ansatte har. HR-direktør i Det Norske Veritas, Gro Gotteberg, nevner i en nylig publisert artikkel at det å innskrenke bruk av hjemmekontor kan oppfattes som å \enquote{[\dots] ta bort noe som er ansett som et gode} \parencite{ingebrigtsen_dnv_2026}. Tilbaketrekking av goder i organisasjonen er til tider nødvendig for å gjøre sentrale forbedringer i selskapet. Argumentene Pål Behrens, leder for arbeidsliv og advokat i Finansforbundet, kommer med, handler i hovedsak om å støtte opp under det sosiale miljøet på arbeidsplassen \parencite{biernat2026}. Utenom dette, er inntrykket at ansatte sitter igjen med følelsen av mangel på visse argumenter, særlig dem som er forankret i bedriftens fordeler med handlingen.

Et sentralt spørsmål blir derfor om fysisk tilstedeværelse er positivt for arbeidsplassen, eller er autonomiteten i hverdagen viktigere for å opprettholde en sunn kultur? Man må ha gode relasjoner til kollegaer for å kunne skape og opprettholde en sunn organisasjonskultur. Dette forebygges ved at ansatte møter opp på arbeidsplassen og skaper rom for andre samtaler enn kun pågående prosjekt. Man ringer sjelden en kollega fra hjemmekontoret for å prate om deres eller eget privatliv.

En innskrenking av muligheten til bruk av hjemmekontor gir rom for å bygge relasjoner. Dersom man mister de små møtene i pausen, kan kvaliteten på tilhørighetsfølelsen hos ansatte få betydelig lavere nivå. Hos nye ansatte vil det også være vanskeligere å bli kjent med kollegaer, og etableringsfasen på arbeidsplassen kan drøyes til lenger enn nødvendig. I en sunn organisasjonskultur er det viktig med gode forhold mellom ansatte. Å stramme inn på hjemmekontor, slik DNB gjorde tidligere i år, kan dermed være med på å styrke de ansattes relasjoner, og arbeidsplassens sosiale miljø kan forbedres.

Tilbøyeligheten med hjemmekontor kan på en annen side skape større overensstemmelse mellom arbeid og privatliv. Dersom man sliter med å få endene til å møtes, kan det gjøre at man har mer å gi av seg selv de dagene en møter opp på jobb. En reaksjon på innstramming av hjemmekontor kan være en form for mistillit til lederen og kan skape følelsen av kontrolltap. Å oppfatte seg selv som et autonomt vesen spiller en sentral rolle i opplevelsen av frihet i hverdagen. Det er ikke slik at innstramming av hjemmekontor direkte fører til tap av autonomitetsfølelsen, men dersom den bakenforliggende begrunnelsen er svak eller mangler forankring i gode argumenter, kan ansatte lettere oppleve det som frihetsinnskrenkende fremfor å se målet bak handlingen. I en sunn organisasjonskultur er det viktig at ansatte har tillit til sin leder og er trygge på at hun tar valg basert på tydelige formål.

Det er også viktig at lederen viser tillit til sine ansatte. Ved bruk av hjemmekontor viser lederen at hun stoler på sine ansatte til å arbeide selvstendig. Problemet hos DNB ligger gjerne i at saken gjelder en innstramming, fremfor en vedvarende situasjon. Som tidligere nevnt kan innstrammingen sees på som fratagelse av en gode. Dersom goden ikke hadde vært der i utgangspunktet, hadde ikke lederen hatt noe utøvelsesrom å minske.

Hos Vend får de ansatte mer mulighet for hjemmekontor av Halvorsen, leder i Vend, likevel om han selv \enquote{trives best rundt kolleger} \parencite{ingebrigtsen_vend_2026}. Ved å gi dem mer frihet og fleksibilitet til å organisere sin egen arbeidsdag, så får tillitt og selvstendighet en større plass i organisasjonskulturen deres. Dette gir dem mulighet for å lage deres egen balanse mellom jobb og fritid, slik at de finner en optimal plan til å gjøre hverdagsoppgaver, alt fra å ta vare barna sine til å delta på fritidsaktiviteter i annen gruppekultur. I DNB vil mange av de ansattes balanser bli omjustert på følge av innstrammingen.

De to bedriftene sine ulike måter å håndtere bruken av hjemmekontor på kan gi uttrykk for forskjellige organisasjonskulturer. Hvor mye de legger vekt på fysisk tilstedeværelse på arbeidsplassen, autonomi og tillitt til ansatte, tyder på ulik grad av artefakter, verdier og normer, og grunnleggende antakelser fra Scheins teori innad i organisasjonene.
